\subsection{RQ4. What is the most frequent type of vulnerability covered in primary studies?}

Software vulnerability detection datasets support different vulnerability types. For example, NVD and SARD benchmark together support 96 types of vulnerabilities. This research question intends to summarize what is the most popular vulnerability types covered by primary studies and what is their frequency. Table \ref{vulnerabilityTypes} shows the statistics regarding the vulnerability types. The column CWE-Type indicates the type of CWE\footnote{\url{https://cwe.mitre.org/}}. There also exists a numerical score for some types of CWEs, which indicate the weakness’ severity. A larger value indicates a higher level of dangerousness and severity\footnote{\url{https://cwe.mitre.org/top25/archive/2021/2021_cwe_top25.html}}. Please note that some frequent types do not have a CWE score, so we denote them as ``-''. There are many categories on the CWE website for vulnerability categorization including \textit{categorization by software development}, \textit{categorization by Hardware design}, and \textit{categorization by research concepts}. The categorization shown in Table \ref{vulnerabilityTypes} is based on \textit{categorization by research concepts} as this categorization is a perfect match for vulnerability types reported in primary studies. 

%\song{put CWE- as the prefix to all the locations you metion the types} \nima{done}
Table \ref{vulnerabilityTypes} indicates that the vulnerability category that receives the highest attendance is related to the improper control of a resource through its lifetime (CWE-664), with a total of 42 studies (18 unique studies). This category primarily involves managing a system's resources, which are created, utilized, and disposed of according to a predefined set of instructions. When a software system fails to follow these guidelines for resource usage, it can lead to unexpected behaviors that create potentially hazardous situations. Attackers can take advantage of these situations to exploit the software system for their own purposes. It is observable that CWE-119 \cite{russell2018automated, liu2020cd, nguyen2020deep, cheng2021deepwukong, yan2021han, fu2022linevul, le2018maximal, cao2022mvd, chen2021neural, lin2019software, filus2020random, li2018vuldeepecker, duan2019vulsniper, wang2020combining} is the most frequent vulnerability type addressed by the primary studies. This vulnerability occurs when a software system attempts to access or write to a memory location outside the permitted boundary of the system's buffer. Attackers can exploit this vulnerability by controlling memory locations and executing their own code or commands, effectively manipulating the system's memory. Although the score for this vulnerability type is not high, the frequency of primary studies addressing it can be a valuable indicator of its significance in terms of detecting and addressing the vulnerability. Within this category, the most severe vulnerability type with a score of 65.93 is CWE-787, which is discussed in 5 primary studies. This vulnerability is considered severe and critical because it can result in the corruption of data, system crashes, or the execution of malicious code. It occurs when a software system attempts to write data beyond the intended buffer, either before the beginning or past the end of the buffer. CWE-22 is another frequent vulnerability type addressed by primary studies accounting for 4 primary studies \cite{fu2022linevul, ghaffarian2021neural, cheng2021deepwukong, cao2022mvd}. This vulnerability is referred to as ``Path Traversal'', where attackers exploit special elements, such as ``..'' or ``/'', to construct their own path and gain unauthorized access to restricted locations. This vulnerability is particularly critical because attackers can use it to modify restricted files or directories in vulnerable software systems, potentially leading to system crashes by uploading malicious code or commands. In critical financial software systems, attackers can even gain access to customers' bank account information. Given the severity of this vulnerability type with a score of 14.69 and the frequency of papers covering it, detecting and addressing this vulnerability is of utmost importance. Therefore, more advanced ML/DL models are needed to effectively detect this type of vulnerability. The least frequent vulnerability type of this family is CWE-120 \cite{russell2018automated, cao2022mvd} which is a classic buffer overflow, accounting for 2 primary studies. This vulnerability occurs when the software attempts to copy a value to the output buffer without first validating its size. If the range of the value is too large for the length of the output buffer, a buffer overflow can occur. While the detection of buffer overflow can be challenging in some cases, static bug detection tools have already addressed this vulnerability, and detection methods are currently available. 

\begin{table}[t]
\caption{Different vulnerability types covered in primary studies.}
\scalebox{0.8}{
\begin{tabular}{p{0.1\linewidth}cccp{0.07\linewidth}l}
\toprule
Category & CWE-Type & Severity Score  & \#Studies & Total & References\\
\midrule
\multirow{1}{*}{CWE-664} & CWE-119 & 5.84 & 14 & \multirow{1}{*}{42(18)} & \cite{russell2018automated, liu2020cd, nguyen2020deep, cheng2021deepwukong, yan2021han, fu2022linevul, le2018maximal, cao2022mvd, chen2021neural, lin2019software, filus2020random, li2018vuldeepecker, duan2019vulsniper, wang2020combining} \\
  %       & CWE-399 & - &8         &       \\
         & CWE-787  & 65.93 & 5         &         &  \cite{wang2020combining, cheng2021deepwukong, fu2022linevul, cao2022mvd, yang2022source} \\
         & CWE-22  & 14.69 &4         &             &  \cite{wang2020combining, cheng2021deepwukong, fu2022linevul, ghaffarian2021neural}  \\
         & CWE-125  & 24.9 & 4         &            &  \cite{wang2020combining, cheng2021deepwukong, cao2022mvd, chen2021neural} \\

         & CWE-400  & - & 4         &             &  \cite{wang2020combining, zou2019mu, cheng2021deepwukong, yang2022source} \\
         & CWE-200  & 4.74 & 3         &            &  \cite{wang2020combining,fu2022linevul, chen2021neural} \\
         & CWE-121  & - & 3         &             &  \cite{russell2018automated, yan2021han, cao2022mvd} \\
         & CWE-122  & - & 3         &             &  \cite{3172102, cao2022mvd, russell2018automated} \\
         & CWE-120  & - & 2         &             &  \cite{russell2018automated, cao2022mvd} \\
         \hline
\multirow{1}{*}{CWE-707}   & CWE-20  & 20.47 & 6 & \multirow{1}{*}{16(9)}               &     \cite{russell2018automated, wang2020combining, cheng2021deepwukong, fu2022linevul, chen2021neural, yang2022source}       \\
                                                 & CWE-78  & 19.55 & 5 &           &    \cite{wang2020combining, cheng2021deepwukong, ghaffarian2021neural, li2021vuldeelocator, yang2022source}      \\
                                                 & CWE-89  & 19.54 & 3 &           &       \cite{ghaffarian2021neural, kronjee2018discovering, wang2020combining}   \\
                                                 & CWE-79  & 46.84 & 2 &           &       \cite{kronjee2018discovering, wang2020combining}   \\
\hline                       
\multirow{1}{*}{CWE-682}    & CWE-190     & 7.12 & 5 & \multirow{1}{*}{5} & \cite{wang2020combining, cheng2021deepwukong, fu2022linevul, zou2019mu, yang2022source} \\
\hline
\multirow{1}{*}{CWE-703} & CWE-476  & 6.54 & 4 & \multirow{1}{*}{4} & \cite{russell2018automated, wang2020combining, cao2022mvd, chen2021neural} \\
    \hline
\multirow{1}{*}{CWE-284} & CWE-284 & - & 2 & \multirow{1}{*}{2} & \cite{aivatoglou2021tree, chen2021neural} \\
    \hline
\multirow{1}{*}{CWE-691} & CWE-362  & - & 2 & \multirow{1}{*}{2} & \cite{chen2021neural, zou2019mu} \\
\hline
\multirow{1}{*}{CWE-1215} & CWE-129 & - & 1 & \multirow{1}{*}{1} & \cite{yang2022source} \\
\hline
\multirow{1}{*}{-} & CWE-789 & - & 1 & \multirow{1}{*}{1} & \cite{yang2022source} \\
\hline
\multirow{1}{*}{SUM} & - & - & - & 73(21) & - \\
\bottomrule
\end{tabular}
}
\label{vulnerabilityTypes}
\end{table}

% CWE78,CWE129,CWE190,CWE400,CWE787,CWE789

\textit{Improper Neutralization - (CWE-707)} 
% \jie{for the first category, can also use the textit to highlight it} \nima{which category? I think it is already highlighted}
Is the second major family of vulnerability types covered by 16 primary studies, including 9 unique primary studies. In this type, the attackers exploit input and output data when they are malformed or not validated properly. There are several scenarios that can lead to data neutralization weaknesses. The first scenario involves checking whether input and output data are safe, while the second scenario involves filtering the input and output data to ensure that any data transformation is done safely. The third scenario involves preventing external attackers from directly manipulating the input and output data, while the fourth scenario relates to the lack of processing input and output data in any circumstances. All of these scenarios can be the root causes of data neutralization weaknesses. As can be seen, CWE-20 is the most frequent type of vulnerability accounting for 6 primary studies. CWE-20 refers to a situation where input validation is not done properly in software systems, making them vulnerable to attacks by malicious individuals who can exploit input data. This occurs when the input data is not verified to be safe or in line with the predefined specifications. The severity and occurrence of this issue are significant, highlighting the need for its detection as it can pose critical risks. CWE-78 is the second major vulnerability type covered by 5 primary studies~\cite{wang2020combining, cheng2021deepwukong, ghaffarian2021neural, li2021vuldeelocator, yang2022source}. This category of security vulnerability pertains to OS command injection, in which an external attacker can construct an OS command by using input data from components that have not been adequately verified. The attacker can then execute harmful commands, potentially causing the system to behave unexpectedly or crash, and putting it in a hazardous state.

\mybox{Answer to RQ4}{red!40}{gray!10}{
(1) The most frequent type of vulnerabilities covered in primary studies is \textit{Improper Control of a Resource Through its Lifetime-(CWE-664)} accounting for 18 unique primary studies. \textit{CWE-119: Improper Restriction of Operations within the Bounds of a Memory Buffer} is the most frequent type of vulnerability in this category accounting for 14 primary studies, and \textit{CWE-787: Out-of-bounds Write} comes subsequently covered by 5 primary studies. The least frequent type of vulnerability is CWE-120 covered by 2 studies.\\
(2) \textit{Improper Neutralization(CWE-707)} is the second major family of vulnerability types covered by 9 unique primary studies in total. In this family, \textit{CWE-20: Improper Input Validation} is the most frequent type covered by 6 primary studies, and \textit{CWE-78: Improper Neutralization of Special Elements used in an OS Command (`OS Command Injection')} comes next with 5 primary studies.\\
(3) Some vulnerability types have a high CVSS score while not sufficiently addressed by existing vulnerability detection studies including but not limited to CWE-79 and CWE-89.
}


